\subsection{p-radical elements}

DEFINITION 1. --- Let K be a field and E an extension of K. An element x of E is said to be p-radical over K if there exists an integer \(m \geqslant 0\) such that \(x^{p^{m}} \in K\) ; the least of these integers is called the height of x (over K).

PROPOSITION 1. --- Let E be an extension of a field K and x a p-radical element of height e over K; put \(a = x^{p^{e}}\) . Then \(a \in K\) and the minimal polynomial of x over K is \(X^{p^{e}} - a\) . Further we have \([\mathbf{K}(x): \mathbf{K}] = p^{e}\) .

It suffices to prove that the polynomial \(X^{p^{e}} - a\) is irreducible in K{[}X{]}. By the definition of height of x we have \(a \notin K^{p}\) , so that the proposition follows from the lemma :

Lemma 1. --- Let K be a field and a an element of K such that \(a \notin \mathbf{K}^p\) . Then for every integer \(e \geqslant 0\) , the polynomial \(f(\mathbf{X}) = \mathbf{X}^{p^e} - a\) is irreducible in \(\mathbf{K}[\mathbf{X}]\) .

Let \(\Omega\) be an algebraic closure of \(\mathbf{K}\) and let \(b\) be the element \(a^{p^{-e}}\) of \(\Omega\) ; we denote by \(g\) the minimal polynomial of \(b\) over \(\mathbf{K}\) . We have \(f(\mathbf{X}) = (\mathbf{X} - b)^{p^e}\) and hence every irreducible polynomial in \(\mathbf{K}[\mathbf{X}]\) which divides \(f\) admits \(b\) as root, and hence is equal to \(g\) . There exists thus (IV, p.~13, Prop. 13) an integer \(q \geqslant 1\) such that \(f = g^q\) ; since \(q\) divides the degree \(p^e\) of \(f\) , there exists an integer \(e'\) such that \(0 \leqslant e' \leqslant e\) and \(q = p^{e'}\) . If \(c\) is the constant term of \(g\) , we have \(-a = c^q\) ; since we assumed that \(a\) does not lie in \(\mathbf{K}^p\) , we must have \(q = 1\) , that is, \(f = g\) , so the lemma follows.

